%----------------------------------------------------------------------------------------
%	PACKAGES AND OTHER DOCUMENT CONFIGURATIONS
%----------------------------------------------------------------------------------------

\documentclass[10pt,stdletter,dateno,sigleft]{newlfm} % Extra options: 'sigleft' for a left-aligned signature, 'stdletternofrom' to remove the from address, 'letterpaper' for US letter paper - consult the newlfm class manual for more options

%\usepackage{charter} % Use the Charter font for the document text
\usepackage{graphics}

\usepackage{standalone}

\usepackage[american]{babel}
\usepackage{csquotes}
\usepackage[backend=biber,style=authoryear,date=year,dashed=false,doi=false,isbn=false,url=false,arxiv=false]{biblatex}
\renewbibmacro{in:}{} % this should go immediately below calling biblatex 
%\DeclareLanguageMapping{american}{american-apa}
\addbibresource{/Users/hectorbahamonde/Bibliografia_PoliSci/library.bib} 


%% with this patch, the newlfm should work (BEGIN)
\usepackage{etoolbox}
\makeatletter
\patchcmd{\@zfancyhead}{\fancy@reset}{\f@nch@reset}{}{}
\patchcmd{\@set@em@up}{\f@ncyolh}{\f@nch@olh}{}{}
\patchcmd{\@set@em@up}{\f@ncyolh}{\f@nch@olh}{}{}
\patchcmd{\@set@em@up}{\f@ncyorh}{\f@nch@orh}{}{}
\makeatother
%% with this patch, the newlfm should work (END)


\newsavebox{\Luiuc}\sbox{\Luiuc}{\parbox[b]{1in}{\vspace{0.35in}
\includegraphics[width=1.2\linewidth]{/Users/hectorbahamonde/Job_Market/logo.png}}} % Company/institution logo at the top left of the page
\makeletterhead{Uiuc}{\Lheader{\usebox{\Luiuc}}}

\newlfmP{sigsize=50pt} % Slightly decrease the height of the signature field
\newlfmP{addrfromphone} % Print a phone number under the sender's address
\newlfmP{addrfromemail} % Print an email address under the sender's address
\PhrPhone{\texttt{p}} % Customize the "Telephone" text
\PhrEmail{\texttt{e}} % Customize the "E-mail" text

\lthUiuc % Print the company/institution logo



% email input begin
\newread\fid
\newcommand{\readfile}[1]% #1 = filename
{\bgroup
  \endlinechar=-1
  \openin\fid=#1
  \read\fid to\filetext
  \loop\ifx\empty\filetext\relax% skip over comments
    \read\fid to\filetext
  \repeat
  \closein\fid
  \global\let\filetext=\filetext
\egroup}
%\readfile{/Users/hectorbahamonde/Bibliografia_PoliSci/email.txt}
% email input end

\usepackage{color} % for colors
\usepackage{hyperref} % for links

\hypersetup{
  colorlinks = true,
  urlcolor = blue,
  pdfpagemode = UseNone
}

%----------------------------------------------------------------------------------------
%	YOUR NAME AND CONTACT INFORMATION
%----------------------------------------------------------------------------------------

\namefrom{\vspace{-3cm}Hector Bahamonde, PhD} % Name

\addrfrom{
{\vspace{-2cm}}\\ % Date
[12pt]
{\bf Hector Bahamonde, PhD}\\
{\color{blue}{\bf Senior Researcher}}\\
\emph{{\color{blue}INVEST Flagship Centre}}\\
{\color{blue}University of Turku}\\
Turku, Finland
}

\phonefrom{+358403621438} % Phone number

\emailfrom{\href{mailto:hector.bahamonde@utu.fi}{hector.bahamonde@utu.fi} \\ 
\texttt{w}: \href{http://www.hectorbahamonde.com}{www.HectorBahamonde.com}\\
\today, download latest version \href{http://github.com/hbahamonde/Job_Market/raw/master/Bahamonde_Cover_Letter.pdf}{\texttt{{\color{red}here}}}}% Link to last version


%----------------------------------------------------------------------------------------
%	ADDRESSEE AND GREETING/CLOSING
%----------------------------------------------------------------------------------------

\greetto{\vspace{-1cm}Dear Members of the Search Committee,} % Greeting text
%\closeline{\vspace{-1cm}\includegraphics[width=3cm]{/Users/hectorbahamonde/Administracion/signature.pdf} \hspace{-4cm}} % Closing text




%----------------------------------------------------------------------------------------

\begin{document}
\begin{newlfm}

%----------------------------------------------------------------------------------------
%	LETTER CONTENT
%----------------------------------------------------------------------------------------

% begining
{\input{/Users/hectorbahamonde/Job_Market/Bahamonde_Cover_letter_P1.tex}\unskip}


In this cover letter I will highlight some details about my research and teaching agendas. As it should become clear, the job ad and my qualifications and experience are a great match. Also, while I am happy in my current position, I look forward to stepping up my career and take bigger responsabilities.

% papers
{\bf Inequality and clientelism.} I kindly invite the reader to read both ``\href{https://journals.sub.uni-hamburg.de/giga/jpla/article/view/1121/1128}{Aiming Right at You: Group versus Individual Clientelistic Targeting in Brazil}'' (Journal of Politics in Latin America, 2018) and ``\href{https://doi.org/10.1057/s41269-020-00174-4}{Still for Sale: The Micro-Dynamics of Vote Selling in the United States, Evidence From a List Experiment}'' ({\input{/Users/hectorbahamonde/research/Vote_Selling/status.txt}\unskip}). In both pieces I challenge the traditional role attributed to {\bf poverty} when explaining clientelism. Instead, I try to switch the focus to {\bf income inequality}. Both papers used top-notch statistical and experimental methods---and when possible, a novel dataset representative at the country level (collected thanks to a generous grant I received). Finally, in ``\href{https://doi.org/10.1016/j.electstud.2022.102497}{Electoral Risk and Vote Buying, Introducing Prospect Theory to the Experimental Study of Clientelism}'' ({\input{/Users/hectorbahamonde/research/Economic_Experiment_Vote_Selling/status.txt}\unskip}) we developed an economic experiment to test the conditions under which parties engage in clientelism. Income inequality plays a big role in explaining vote buying.

{\bf Economic inequality and democracy.} We have published a paper that explores the relationship between state capacity, democracy and inequality. Using time-series econometrics we find in ``\href{https://doi.org/10.1016/j.ejpoleco.2021.102048}{Inclusive Institutions, Unequal Outcomes: Democracy, State Capacity, and Income Inequality}'' ({\input{/Users/hectorbahamonde/research/Inequality_State_Capacity/status.txt}\unskip}) that democratic rule combined with high state infrastructural power produce \emph{higher} levels of income inequality over time. This relationship operates through the positive effect of high-capacity democratic context on investor confidence, FDIs and financial development. Within the same line of research, in ``\href{https://doi.org/10.1111/pops.12940}{\fullcite{Bahamonde:2023}}'' we discover that female candidates that \emph{look like} they have working-class occupations are electorally penalized when running for elections--male candidates are not affected by their working-class looks.

%{\bf State building.} Using a novel dataset on historical earthquakes to proxy state capacity, in ``\href{https://github.com/hbahamonde/Earthquake_Paper/raw/master/Bahamonde_Earthquake_Paper.pdf}{{\input{/Users/hectorbahamonde/research/Earthquake_Paper/title.txt}\unskip}}'' ({\input{/Users/hectorbahamonde/research/Earthquake_Paper/status.txt}\unskip}) I explain that sectoral conflicts between the landed and industrial elites fostered inter-elite compromises that lead to higher levels of state capacity over time. The paper seeks to analyze how economic structural transformations in Latin America helped states to make institutional investments that lead to the formation of states with higher capacity. I leverage fine-grained historical case study comparisons, sectoral economic outputs from 1900 to the present, Bayesian time-series econometric techniques, hazard models, and a novel earthquake dataset that covers sub-national death tolls from 1900 to the present to measure state capacity. This project builds on the fiscal sociology literature and the dual-sector economy model. 


% agenda: structure / behavior
As a comparativist and political economist, I believe that advanced methods should be used to answer big questions. In this sense, my research also has a disciplinary agenda. My scholarly work seeks to examine classic problems in comparative political development (e.g. state capacity, inequality, clientelism, etc.), while at the same time incorporating cutting-edge econometric and experimental techniques. Similarly, my working papers and experiments, are concerned with fundamental questions regarding democratic theory.

%{\bf Digitalization: a research agenda.} Cities are constantly producing digital data. From MetroCard transactions, social media posts to Google searches. I feel particularly intrigued by new ways of exploiting such data so we can answer questions about inequality and clientelism. During my three years at WZB (IPI unit), I intend to move forward this agenda, particularly by publishing several papers that I have in the pipeline. First, I have two papers using conjoint data on clientelism (please check my \href{http://github.com/hbahamonde/Job_Market/raw/master/Bahamonde_CV.pdf}{\texttt{CV}}). All are at different stages of the publication process. Those digital data were collected via ICTs in an experimental way. Second, I have collected a digitalized population dataset on daily public transportation and COVID contagion in Santiago de Chile, the capital city of one of the most unequal countries in the world. I intend to study how low-income municipalities systematically bore the cost of the COVID pandemic. This is a really interesting project as it shows that (1) contagion thresholds to restrict mobility in low-income municipalities were higher and that (2) lockdown policies during daytime were systematically ignored by local authorities but heavily enforced during nighttime---which is curfew time, often used to protest against the government. These findings are relevant as they go in line with the recent ``politics of weakness'' literature. I use hazard models, synthetic control and regression discontinuity design methods. To conclude this portion of the statement, my research agenda on digitalization not only involves publishing these papers, but also organizing talks and presenting at conferences, as well as collaborating with other WZB affiliates.

% teaching (topics) // service
As an {\bf instructor}, I am interested in courses that are carefully designed to answer big questions in comparative politics, not only from a Latin American perspective, but also across the broader spectrum. As a political economist, I mostly focus on economic development, economic history, institutional development and democratic development/underdevelopment. I am also interested in political participation, competitive authoritarianisms, welfare politics, party politics and other relevant topics. I would also feel very comfortable teaching both basic and advanced {\bf methods} courses. At Rutgers, Tulane and Turku, I always received excellent teaching evaluations. Please, send me an email to receive the most recent one, or let me know if you need teaching references. {\bf Service} to the department and its extended programs is always a top priority. I firmly believe that being part of a teaching/research community requires active involvement. I am looking forward to do my part in this regard. {\bf Collaborative research with my students} is another priority. For example, as an undergraduate student I worked very close with faculty members doing research. I learned so much outside of the classroom by doing research that it would \emph{not} be natural for me to not replicate my experience with my own students.



% ending
{More information, \href{http://www.hectorbahamonde.com/teaching/}{\texttt{syllabi}}, my \href{http://github.com/hbahamonde/Job_Market/raw/master/Bahamonde_Research_Statement.pdf}{\texttt{research}}, \href{http://github.com/hbahamonde/Job_Market/raw/master/Bahamonde_Teaching_Statement.pdf}{\texttt{teaching}} and \href{http://github.com/hbahamonde/Job_Market/raw/master/Bahamonde_Diversity_Statement.pdf}{\texttt{diversity}} statements, as well as other \href{http://www.hectorbahamonde.com/research/}{\texttt{papers}} are available on my website: \href{http://www.hectorbahamonde.com}{\texttt{www.HectorBahamonde.com}}. Thank you for considering my application. I look forward to hearing from you.\unskip}


%----------------------------------------------------------------------------------------

\end{newlfm}
\end{document}
